\begin{figure}[htbp]
\centering
\begin{tikzpicture}[causal,x=1.5cm,y=1.4cm]
\draw[ed,thin] (0,0)--(5.5,0) node[above]{$x$};
\draw[ed,thin] (0,0)--(0,1.45) node[left]{$y$};
\foreach \x in {1,2,3,4}{
 \draw[thin] (\x,-.03)--(\x,.03);
 \node[below] at (\x,-.04){$\x$};}
\draw[thin] (-.04,1)--(.04,1);
\node[left] at (-.03,1){$1$};
\node[left] at (-.03,0){$0$};
\foreach \x/\y in {1/0,2/0,3/1,4/0}{\fill (\x,\y) circle (1.4pt);}
\node[below left] at (2.8,.95){$f_3$};
\begin{scope}[yshift=-2.9cm]
\draw[ed,thin] (0,0)--(5.5,0) node[above]{$x$};
\node[left] at (0,.42){$P_X$};
% 相对柱高取自原 PDF 矢量：31.8901、31.8901、10.6300、42.5202 点。
\fill[plotgray] (.95,0) rectangle (1.05,.6);
\fill[plotgray] (1.95,0) rectangle (2.05,.6);
\fill[plotgray] (2.95,0) rectangle (3.05,.2);
\fill[plotgray] (3.95,0) rectangle (4.05,.8);
\end{scope}
\end{tikzpicture}
\caption{在这个例子中，即使沿因果方向预测，半监督学习也能降低损失。由于对每个 $x$，标签零的先验可能性都高于标签一，因此，选择一个仅在 $p(x)$ 最小的点 $x$ 处取值为一的函数，就能使错误数的期望最小。此处，该点为 $x=3$。}
\label{fig:ch05-ssl-counterexample}
\end{figure}
